\section{Dynamic Programming}
\subsection{SOS DP (Sum Over Subsets)}
Complexity $O(B \cdot 2^B)$. Used for gcd/lcm/OR/AND counting. \texttt{zeta\_subset}: for each mask, sum over its submasks. \texttt{zeta\_superset}: for each mask, sum over its supermasks.
\begin{minted}{c++}
void zeta_subset(vector<int64_t> &dp, int bits) {
  for (int b = 0; b < bits; b++)
    for (int mask = 0; mask < (1LL << bits); mask++)
      if (mask >> b & 1)
        dp[mask] += dp[mask ^ (1LL << b)];
}
void zeta_superset(vector<int64_t> &dp, int bits) {
  for (int b = 0; b < bits; b++)
    for (int mask = 0; mask < (1LL << bits); mask++)
      if (!(mask >> b & 1))
        dp[mask] += dp[mask | (1LL << b)];
}
\end{minted}
\begin{minted}{c++}
void Solve() {
  ll n; cin >> n;
  vector<ll> arr(n);
  for (auto& x : arr) cin >> x;
  ll N = 1 + __lg(*max_element(arr.begin(),
  arr.end()));
  vector<ll> sos(1LL << N, 0), sos2(1LL << N, 0);
  for (auto it : arr) {
    sos[it]++;
    sos2[((1LL << N) - 1) ^ it]++;
  }
  for (int i = 0; i < N; ++i)
    for (int mask = 0; mask < (1LL << N); ++mask)
      if (mask & (1LL << i))
        sos[mask] += sos[mask ^ (1LL << i)],
        sos2[mask] += sos2[mask ^ (1LL << i)];
  for (auto& it : arr) {
    // count y: x|y == x
    cout << sos[it] << ' ';
    // count y: x&y == x
    cout << sos2[((1LL << N) - 1) ^ it] << ' ';
    // count y: x&y != 0
    cout << n - sos[((1LL << N) - 1) ^ it] << ' ';
  }
}
\end{minted}

\subsection{Knuth Optimization}
For DP on intervals where optimal split point $p[i][j]$ is monotone in both $i$ and $j$:
$dp[i][j] = \min_{i<k<j} (dp[i][k] + dp[k][j]) + f(i, j)$
Reduces $O(N^3)$ to $O(N^2)$.
\textbf{Note:} Speeds up interval DP from $O(N^3)$ to $O(N^2)$ if costs are monotonic.
\begin{minted}{c++}
for (int i = 0; i <= n; i++) {
  p[i][i] = i, dp[i][i] = 0;
  if (i < n) dp[i][i+1] = cost(i, i+1), p[i][i+1] = i;
}
for (int len = 2; len <= n; len++) {
  for (int i = 0; i + len <= n; i++) {
    int j = i + len;
    dp[i][j] = LLONG_MAX;
    for (int k = p[i][j-1]; k <= p[i+1][j]; k++) {
      ll val = dp[i][k] + dp[k][j] + cost(i, j);
      if (val < dp[i][j]) {
        dp[i][j] = val;
        p[i][j] = k;
      }
    }
  }
}
\end{minted}
\subsection{Max Rectangular Sum in Matrix}
Key idea: Fix top and bottom rows, collapse 2D into 1D column sums, then apply Kadane's algorithm for max subarray. $O(N^2 M)$.
\begin{minted}{c++}
int64_t maxRectSum(vector<vector<int>>& a) {
  int n = a.size(), m = a[0].size();
  int64_t ans = -4e18; // Safe negative infinity
  for (int top = 0; top < n; top++) {
    vector<int64_t> col(m, 0);
    for (int bot = top; bot < n; bot++) {
      for (int j = 0; j < m; j++)
        col[j] += a[bot][j];
      int64_t cur = 0;
      for (int j = 0; j < m; j++) {
        cur = max(col[j], cur + col[j]);
        ans = max(ans, cur);
      }
    }
  }
  return ans;
}
\end{minted}
\subsection{Subset Sum (Bitset)}
Works when sum of values $\le N$. $O(N \sqrt{N} / 64)$.
Uses bitset DP with binary splitting optimization.
\texttt{dp.test(k)} is true if sum $k$ can be formed.
\begin{minted}{c++}
const int N = 1e6 + 6;
bitset<N> dp;
dp.set(0);
for (auto &[w, c] : comp) {
  for (int i = 0; (1LL << i) <= c; ++i) {
    dp |= dp << (w * (1LL << i));
    c -= 1LL << i;
  }
  dp |= dp << (w * c);
}
dp.test(k) // if sum k is possible
\end{minted}
OR convolution can be done using SOS DP.
\subsection{Exclusion DP (GCD/LCM counting)}
Count pairs/k-tuples with given gcd. Let $d[i]$ = count of elements divisible by $i$. Then: \\
\textbf{Ordered pairs} $(a,b)$, $(a,b) \neq (b,a)$: $f[i] = d[i] \cdot d[i]$ \\
\textbf{Unordered (with rep)}: $f[i] = d[i](d[i]+1)/2$ \\
\textbf{Unordered (no rep)}: $f[i] = d[i](d[i]-1)/2$ \\
\textbf{Ordered $k$-tuples}: $f[i] = (d[i])^k$ \\
\textbf{Unordered (no rep)}: $f[i] = \binom{d[i]}{k}$ \\
\textbf{Unordered (with rep)}: $f[i] = \binom{d[i]+k-1}{k}$


\subsection{Li Chao Tree (CHT)}
Maintains lines $y = mx + b$ and segments. Answers min $y$-value at query $x$. Line add/query $O(\log C)$, segment add $O(\log^2 C)$ ($C$ = coordinate range).

When to use: DP transitions $dp[j] = \min_i (m_i \cdot x_j + b_i)$ without monotonic slopes/queries. Nodes allocated dynamically. For max queries: insert $\{-m, -b\}$ and negate answer.
\textbf{Note:} Works over dynamic $[-10^9, 10^9]$ range without coordinate compression.
\begin{minted}{c++}
struct LiChao {
  struct Line {
    ll m, b;
    ll eval(ll x) const { return m * x + b; }
  };
  static constexpr ll INF = 2e18;
  struct Node {
    int l = -1, r = -1;
    Line line = {0, INF};
  };
  vector<Node> t = {Node()};
  void add_line(Line nw, int u = 0, ll l = -1e9, ll r = 1e9) {
    ll mid = l + (r - l) / 2;
    bool bl = nw.eval(l) < t[u].line.eval(l);
    bool bm = nw.eval(mid) < t[u].line.eval(mid);
    if (bm) swap(t[u].line, nw);
    if (l == r) return;
    if (bl != bm) {
      if (t[u].l == -1) { t[u].l = t.size(); t.emplace_back(); }
      add_line(nw, t[u].l, l, mid);
    } else {
      if (t[u].r == -1) { t[u].r = t.size(); t.emplace_back(); }
      add_line(nw, t[u].r, mid + 1, r);
    }
  }
  void add_segment(ll ql, ll qr, Line nw, int u = 0, ll l = -1e9, ll r = 1e9) {
    if (qr < l || r < ql) return;
    if (ql <= l && r <= qr) { add_line(nw, u, l, r); return; }
    ll mid = l + (r - l) / 2;
    if (ql <= mid) {
      if (t[u].l == -1) { t[u].l = t.size(); t.emplace_back(); }
      add_segment(ql, qr, nw, t[u].l, l, mid);
    }
    if (qr > mid) {
      if (t[u].r == -1) { t[u].r = t.size(); t.emplace_back(); }
      add_segment(ql, qr, nw, t[u].r, mid + 1, r);
    }
  }
  ll query(ll x, int u = 0, ll l = -1e9, ll r = 1e9) {
    if (u == -1) return INF;
    ll res = t[u].line.eval(x);
    if (l == r) return res;
    ll mid = l + (r - l) / 2;
    if (x <= mid) return min(res, query(x, t[u].l, l, mid));
    return min(res, query(x, t[u].r, mid + 1, r));
  }
};
// For max queries: insert {-m, -b} and answer is -query(x)
\end{minted}
